\section{Replacing the end and retaining the area}\label{sec:introduction}

We ask whether a Cha--Khuri--Sakovich hyperboloidal end can be replaced
by an asymptotically flat end while leaving a prescribed compact
interior unchanged and making the ADM mass approach the Lorentz norm
of the hyperboloidal charge.
For the Penrose inequality, a mass comparison alone is insufficient.
A surface enclosing the inner boundary can travel through the modified
end, so its area must remain controlled wherever the replacement occurs.

We construct replacements that preserve the dominant energy condition
and give a lower metric comparison on the entire exterior. This controls
every enclosing surface at once, including disconnected surfaces and
surfaces touching the original boundary. The construction applies to
the precise asymptotic class of Cha, Khuri and Sakovich
\cite{CKS}. The end replacement is the main geometric result; the sharp
Bondi inequality follows by applying the numerical asymptotically flat
inequality to each completed replacement.

\subsection{The exterior, its charges and its enclosing cuts}
Let $\Omega$ be a connected orientable smooth three-manifold with
nonempty compact smooth boundary $S$. Let $g$ be a smooth Riemannian
metric and $K$ a smooth symmetric covariant two-tensor, both smooth up
to $S$. We require completeness as a metric space with $S$ included,
and exactly one coordinate end with compact complement.
With $G=c=1$ and zero cosmological constant, the constraint densities are
\begin{equation}\label{eq:constraints-definition}
 16\pi\mu=R_g+(\tr_gK)^2-|K|_g^2,
 \qquad 8\pi J=\diver_g\bigl(K-(\tr_gK)g\bigr).
\end{equation}
The dominant energy condition (DEC) is $\mu\geq |J|_g$.
For a two-sided surface with normal $\nu$ pointing into its end side,
write
\[
 H=\diver_{\mathrm{surface}}\nu,\qquad
 \theta_+=H+\tr_{\mathrm{surface}}K,\qquad
 \theta_-=H-\tr_{\mathrm{surface}}K.
\]
Thus a future marginally outer trapped surface has $\theta_+=0$.
Our spacetime convention is
$K(Y,Z)=\overline g(\overline\nabla_Yn_{\mathrm{future}},Z)$;
the future unit hyperboloid in Minkowski space has $K=g$.

An \emph{enclosing cut} is $\Gamma=\partial D$, the intrinsic manifold
boundary of a connected smooth codimension-zero submanifold-with-boundary
$D\subset\Omega$ that is closed in $\Omega$, contains the whole
sufficiently distant end, and has manifold interior in
$\operatorname{int}\Omega$. Its full boundary must be compact, smooth,
embedded and two-sided. Every component and every portion coinciding
with $S$ contributes its area. A cut need not be connected, trapped or
minimal. Set
\begin{equation}\label{eq:area-definition}
 A_{\min,g}(S)=\inf_\Gamma\operatorname{Area}_g(\Gamma).
\end{equation}
Taking $D=\Omega$ shows that the class includes $S$ itself. We neither
assume attainment of the infimum nor identify it with
$\operatorname{Area}_g(S)$. This distinction matters for the Penrose
inequality: Ben-Dov's spherical example violates the bound using the
apparent horizon's own area, while leaving the minimum-enclosing-area
formulation intact \cite[Sections~1 and~4]{BenDov}.

Let $\sigma$ be the unit round metric on $\Sph$. In coordinates
$(r,\omega)$ on the end, put
\[
 h=(1+r^2)^{-1}\mathrm dr^2+r^2\sigma,\qquad g=h+a,\qquad K=h+b.
\]
The CKS class in this paper consists exactly of the expansions
\begin{equation}\label{eq:cks}
\begin{aligned}
 a_{rr}&=r^{-5}m^r+O_3(r^{-6}),&a_{rA}&=O_3(r^{-3}),&
 a_{AB}&=r^{-1}m^g_{AB}+O_3(r^{-2}),\\
 b_{rr}&=O_2(r^{-5}),&b_{rA}&=O_2(r^{-3}),&
 b_{AB}&=r^{-1}m^K_{AB}+O_2(r^{-2}).
\end{aligned}
\end{equation}
The leading angular fields are smooth. For a coordinate component,
$O_j(r^{-q})$ bounds $r^{q+i}\partial_r^i\partial_\omega^\alpha$
through $i+|\alpha|\leq j$ in a finite sphere atlas. Equivalently,
one may use angular derivatives and scaled radial derivatives
$r\partial_r$. No smallness, pointwise sign or symmetry of the leading
fields is required. These component rates and derivative orders are
those of \cite[Equations~(2.1)--(2.3)]{CKS}.

Writing $n=(n^1,n^2,n^3)$ for the inclusion $\Sph\subset\R^3$,
define
\begin{equation}\label{eq:bondi}
 \begin{gathered}
 M=\tr_\sigma(m^g+2m^K)+2m^r,\qquad
 (E_B,P_B)=\frac1{16\pi}\int_{\Sph}(1,n)M\,\mathrm dA_\sigma,\\
 m_B=\sqrt{E_B^2-|P_B|^2}.
 \end{gathered}
\end{equation}
We call $(E_B,P_B)$ the \emph{Bondi charge of the initial data}.
All statements below concern the explicit CKS functional
\eqref{eq:bondi} and its Lorentz norm; no spacetime development is
assumed. Identifying an initial-data charge with the Bondi--Sachs charge
of a cut of null infinity requires an appropriate compatible
development and asymptotics; see, for example,
\cite[Theorem~5.3]{CJL}. We assume $E_B>|P_B|$.
A chart is \emph{balanced} if $P_B=0$;
then $E_B=m_B>0$. The second tensor is essential to the aspect $M$.
For comparison with Chen, Wang and Yau, set $p=K-g$. Its leading
tangential coefficient is $m^K-m^g$, so their mass aspect is
\[
 \tfrac32\tr_\sigma m^g+\tr_\sigma(m^K-m^g)+m^r=M/2.
\]
Their charge factor $1/(8\pi)$ therefore gives exactly
\eqref{eq:bondi} \cite[Definitions~1.3--1.4]{CWY}.
Chen, Wang and Yau obtain rest-frame foliations in their asymptotic
setting \cite[Corollary~4.3]{CWY}. Lemma~\ref{lem:balance} proves the
Lorentz covariance and preservation of the finite derivative class
\eqref{eq:cks} needed here.

\subsection{The geometric theorem and its numerical consequence}
For an asymptotically flat end with Cartesian coordinates $x=rn$,
our ADM conventions are
\begin{align}
 E&=\frac1{16\pi}\lim_{r\to\infty}\int_{|x|=r}
 (\partial_jg_{ij}-\partial_ig_{jj})n^i\,\mathrm dA_\delta,
 \label{eq:adm-energy}\\
 P_i&=\frac1{8\pi}\lim_{r\to\infty}\int_{|x|=r}
 \bigl(K_{ij}-(\tr_gK)g_{ij}\bigr)n^j\,\mathrm dA_\delta.
 \label{eq:adm-momentum}
\end{align}
In Cartesian coordinates $O_j(r^{-q})$ has the usual differentiated
decay: an order-$i$ derivative is $O(r^{-q-i})$.

\begin{theorem}[Area-controlled end replacement]\label{prop:replacement}
Let $(\Omega,g,K)$ be a smooth connected orientable three-dimensional
exterior, complete including its nonempty compact smooth boundary,
with exactly one end of class \eqref{eq:cks}. Suppose DEC holds and
$E_B>|P_B|$. There is a balanced chart such that, for every sufficiently
large $R$, smooth data $(G_R,k_R)$ on the same exterior have the following
properties:
\begin{enumerate}
\item They agree with $(g,K)$ on the compact side of $r=R$ and satisfy DEC.
\item For numbers $0\leq\varepsilon_R<1$ with $\varepsilon_R\to0$,
      $G_R\geq(1-\varepsilon_R)g$ everywhere on $\Omega$.
\item The end is asymptotically flat, with $G_R-\delta=O_j(r^{-1})$
      for every fixed finite $j$, and $k_R=0$ sufficiently far out.
      The constraints are integrable and its charges satisfy
      $P_R=0$, $E_R=m_B+\eta_R$, where $0\leq\eta_R\leq2R^{-1/2}$.
\end{enumerate}
The new exterior is complete and
$A_{\min,G_R}(S)\geq(1-\varepsilon_R)A_{\min,g}(S)$.
No condition on the boundary expansion is needed.
\end{theorem}

The theorem separates the geometry of the end from a numerical
inequality. In particular it retains the boundary's original null
expansion, whatever its sign. For a weakly future trapped boundary,
the numerical spacetime Penrose inequality gives the following result.

\begin{corollary}[Sharp Bondi inequality in the CKS class]
\label{cor:cks-inequality}
Under the hypotheses of Theorem~\ref{prop:replacement}, assume
$\theta_+(S)\leq0$, with normal into $\Omega$, and
$A_{\min,g}(S)>0$. Then
\begin{equation}\label{eq:main}
 \sqrt{E_B^2-|P_B|^2}\geq
 \sqrt{\frac{A_{\min,g}(S)}{16\pi}}.
\end{equation}
The boundary may be disconnected. For every $m>0$ there is a
horizon-regular Schwarzschild hyperboloidal exterior attaining equality
with $m_B=m$ and $A_{\min,g}(S)=16\pi m^2$.
\end{corollary}

This gives the sharp Bondi--Penrose bound in the specified CKS class,
using the numerical spacetime Penrose theorem of
\cite[Definitions~1.1--1.2 and Theorem~1.3]{NumericalPenrose} in
dimension three. Section~\ref{sec:companion-input} states that input
and checks its full hypotheses for each fixed replacement. The
geometric theorem also gives a conditional implication from the
corresponding asymptotically flat inequality, recorded in
Section~\ref{sec:conditional}. The sharpness examples exclude additional
closed horizons of either time orientation.

\subsection{Predecessors and the construction}
Penrose's mass--area comparison arose from gravitational collapse,
cosmic censorship and the area theorem \cite{Penrose}.
For time-symmetric asymptotically flat data, Huisken and Ilmanen's
weak inverse mean curvature flow proves the bound for each connected
component of an outermost minimal horizon \cite{HI}. Bray's conformal
flow gives the sharp bound for the total area of a possibly disconnected
outer-minimizing minimal horizon \cite{Bray}. The extension to a general
second fundamental form must retain an area comparison while controlling
both constraint densities.

In the hyperboloidal setting, Cha, Khuri and Sakovich reduce the sharp
inequality to a generalized Jang equation coupled to weak inverse mean
curvature flow, with the required boundary behavior and asymptotics and
a connected outermost apparent horizon \cite[Theorem~3.2]{CKS}.
Their prescribed-warping existence theorem solves the generalized Jang
equation alone \cite[Theorem~3.4]{CKS}; it does not establish existence
for the coupled system. Sakovich's ordinary Jang argument proves
positive energy in a weighted asymptotically hyperbolic class
\cite[Theorem~1.1]{Sakovich}.

Null-geometric approaches make the role of the limiting cut and observer
explicit. Mars and Soria compare a trapped cross-section with the
limiting Hawking energy along a distinguished null foliation;
identifying that limit with Bondi energy requires their large-sphere
condition \cite[Theorems~2--3]{MarsSoria}. In a controlled perturbative
vacuum regime near Schwarzschild, Alexakis constructs a nearby
marginally trapped sphere and an asymptotically round null foliation
giving a Bondi-energy inequality at its new limiting cut
\cite[Theorem~1.3 and Section~2.1]{Alexakis}.

For nonsymmetric initial data, Allen, Bryden, Kazaras and Khuri prove
a universal suboptimal-constant inequality for invariant ADM mass in
the asymptotically flat case and for the chosen end's energy under
Wang-type hyperboloidal asymptotics
\cite[Theorems~1.1 and~1.3]{ABKK}. Their area is the least enclosing
area of a selected outermost apparent horizon. Ellithy's
future-development approach bounds actual horizon-section area under
quasi-final-state and further Bondi--Sachs hypotheses
\cite[Theorem~4.12, Remark~4.14 and Theorem~C.1]{Ellithy}.
Its final Bondi quantity is defined in the specified Bondi--Sachs
foliation. These results distinguish the initial-data norm, a
foliation's energy, and the area being bounded.

Our construction is related to earlier geometric changes of an end.
The radial-lapse ansatz belongs to Bartnik's quasi-spherical framework
\cite{Bartnik}\cite[Section~5]{BartnikICM}.
Chru\'sciel and Delay bend an exactly hyperbolic end to a flat end
through a Minkowski graph, preserving the compact side and DEC
\cite[proof of Theorem~4.2]{ChruscielDelay}. Their model has zero charge
and uses the opposite time orientation. Corvino and Huang's localized
deformation theory addresses the additional difficulty that a metric
change also changes the norm of the momentum density
\cite[Introduction and Equation~(1.1)]{CorvinoHuang}.
Ellithy's late-time comparison likewise attaches a quasi-spherical
Riemannian end before applying inverse mean curvature flow
\cite[Appendix~C, proof of Theorem~C.1]{Ellithy}.
The target here is a DEC replacement of the fixed CKS initial-data
exterior, with positive charge approximated by ADM mass and a global
lower metric bound protecting every enclosing cut. The collar and
bending estimates establish these properties directly.

The proof begins with a radial vacuum model in
Section~\ref{sec:model}. It shows how the tangential tensor trace can
be reduced to zero while retaining a mass variable. An arbitrary angular
mass aspect creates two additional terms: a Laplacian in energy and a
gradient in momentum. Heat evolution cancels the first, and a trace-free
sphere tensor cancels the second. Its divergence equation has precisely
a first-harmonic obstruction; Lorentz balancing removes that obstruction.

Before bending, the original leaves must be made round. Because the
input DEC may be saturated, this cannot be justified by smallness alone.
Section~\ref{sec:geometry} supplies the constraint identities, and
Section~\ref{sec:collar} constructs a collar whose added constraint
vector lies in the future causal cone. The estimates retain the
activation cutoff as a factor, even at its flat initial endpoint.
Section~\ref{sec:bending} then performs the angularly corrected bending.
Heat starts while the tangential tensor trace still has its
hyperboloidal value; the trace-free correction enters as bending begins.
Section~\ref{sec:transfer} computes the ADM charges at fixed $R$,
compares every cut, and applies the numerical theorem before letting
$R\to\infty$. Section~\ref{sec:sharpness} supplies actual equality
examples. Vanishing comparison errors alone do not classify equality
for general hyperboloidal data.
